\speciestitle{Temperature}{T}{Temperature}{%
\item[Useful range:] 261\,--\,0.001\,hPa
\item[Contact:] Michael J. Schwartz,
  \email{Michael.J.Schwartz@jpl.nasa.gov}}
\label{sec:StandardTemperature}%
\subsection*{Introduction}

The MLS \thisv temperature product is similar to the v2.2
product that is described in~\citet{SchwartzEtal08}.  MLS temperature
is retrieved from bands near O\cs2 spectral lines at
118\,GHz and 239\,GHz that are measured with MLS radiometers R1A/B and R3,
respectively.  
The isotopic 239-GHz line is the primary source of
temperature information in the troposphere, while the 118-GHz line is
the primary source of temperature in the stratosphere and above.  
MLS \thisv temperature has a \texttildelow\,$-$1\,K  bias with
respect to correlative measurements in the troposphere
and stratosphere, with 2\,--\,3\,K peak-to-peak additional systematic
vertical structure.  Table~\ref{tab:TSummary} summarizes the measurement
precision, resolution, and modeled and observed biases.  The following
sections provide details.

\subsection*{Differences between \thisv and v2.2 }
The MLS \thisv temperature retrieval algorithms are largely unchanged from
those of v2.2, using the same subsets of the same radiance bands, so the
resulting products are very similar. An exception is the \mbox{316-hPa}
level, which in \thisv is noisier and has larger biases relative to
analyses than in v2.2.  The \thisv temperature \mbox{316-hPa} level is not
recommended for scientific use.  Version \thisv has eight more retrieval
levels in the upper stratosphere, giving 12 levels per decade from the
surface to 1\,hPa.  Noise and biases have been reduced at ``chunk
boundaries'', the breaks between the 10-profile blocks of data that are
concurrently retrieved by the 2-D algorithms.  The non-convergence of
the retrieval over substantial sections of the polar autumn in v2.2 has
been eliminated in version \thisv.  Vertical smoothing has been reduced in
the mesosphere and lower thermosphere, improving vertical resolution at
a cost of less than a factor of two in precision, but also resulting in
what, in preliminary validation, appears to be some vertically
oscillating profiles in the mesopause region, particularly at the
equator. The \emph{a priori} temperature profiles used in \thisv are
consistently GEOS-5.2, while v2.2 used GEOS-5.1 before September of
2008.  GEOS-5.1 has a low bias in the upper stratosphere of
0\,--\,15\,K, particularly at high latitudes, and resulted in biases of
\texttildelow1\,K near the stratopause.

\subsection*{Resolution}
\onestandardavk{Temperature_HR}{Temperature}%
The vertical and horizontal resolution of the MLS temperature
measurement is shown by averaging kernels in
Figure~\ref{fig:AvkTemperature_HR}.  Vertical resolution, shown on the
left panel, is \texttildelow5\,km from 261\,hPa to 100\,hPa, improves to
3.6\,km at 31.6\,hPa and then degrades to 4.3\,km at 10\,hPa, 5.5\,km at
3.16\,hPa and 6\,km at 0.01\,hPa.  Along track resolution is
\texttildelow170\,km from 261\,hPa to 0.1\,hPa and degrades to 220\,km at
0.001\,hPa.  The cross-track resolution is set by the 6-km width of the
MLS 240-GHz field of view in the troposphere and by the 12-km width of
the MLS 118-GHz field of view in the stratosphere and above.  The
longitudinal separation of MLS measurements from a given day, which is
determined by the Aura orbit, is 10\degsym--\,20\degsym\ over middle and
low latitudes and much finer in polar regions.

\subsection*{Precision}

The precision of the MLS \thisv temperature measurement is summarized in
Table~\ref{tab:TSummary}.  Precision is the random component of
measurements which would average down if the measurement were repeated.
The retrieval software returns an estimate of precision based upon the
propagation of radiometric noise and \emph{a~priori} uncertainties through the
measurement system.  These values, which range from 0.6\,K in the lower
stratosphere to 2.5\,K in the mesosphere, are given, for selected
levels, in column~2.  Column~3 gives the rms of
differences of values from successive orbits (divided by the square-root
of two as we are looking at the difference of two noisy signals) for
latitudes and seasons where longitudinal variability is small and/or is
a function only of local solar time.  The smallest values found, which
are for high-latitude summer, are taken to be those least impacted by
atmospheric variability, and are what is reported in column~3.  These
values are slightly larger than those estimated by the measurement
system in the troposphere and lower stratosphere, and a factor of
\texttildelow1.4 larger from the middle stratosphere through the mesosphere.

\subsection*{Accuracy}

A substantial study of sources of systematic error in MLS measurements
was done as a part of the validation of v2.2, and as the measurement
system is substantially unchanged, those results are repeated here.  The
accuracy of the v2.2 temperature measurements was estimated both by
modeling the impact of uncertainties in measurement and retrieval
parameters that could lead to systematic errors, and through comparisons
with correlative data sets.  Column~5 of Table~\ref{tab:TSummary} gives
estimates from the propagation of parameter uncertainties, as discussed
in~\citet{SchwartzEtal08}.  This estimate is broken into two pieces.
The first term was modeled as amplifier non-linearity, referred to as
``gain compression,'' and was believed to have a known sign, as gain is
known to drop at high background signal levels. Correction of these
linearity's was a goal of \thisv, but closer examination of the simple
non-linearity model found that it did not close foreword model and
measured radiances as expected.  It had been hoped that better radiance
closure would permit the use of more radiances in the middle of the
118-GHz O$_2$ band, giving better resolution, precision and accuracy in
the upper stratosphere and better accuracy everywhere.  This work is
still ongoing, and it is hoped that advances will manifest in
improvements in a future version.

The second term of column~5 combines 2-$\sigma$ estimates of other
sources of systematic uncertainty, such as spectroscopic parameters,
retrieval numerics and pointing, for which the sign of resulting bias is
unknown.  Gain compression terms range from $-$1.5\,K~to~$+$4.5\,K, and
predicted vertical structure is very similar to observed biases relative
to correlative data in the troposphere and lower stratosphere.  The
terms of unknown sign are of \texttildelow2\,K magnitude over most of the
retrieval range, increasing to 5\,K at 261\,hPa and to 3\,K at
0.001\,hPa.


Column~6 contains estimates of bias based upon comparisons with analyses
and with other previously-validated satellite-based measurements.  In
the troposphere and lower stratosphere, the observed biases between
MLS and most correlative data sets are
consistent to within \texttildelow1.5\,K, and have vertical oscillation with 
an amplitude of 2\,--\,3\,K and a vertical frequency of
about 1.5~cycles per decade of pressure.  A global average of
correlative measurements is shown in Figure~\ref{fig:TempSummary2}.

\subsection*{Data screening}
\begin{description}
\item[Pressure range:] \screeningrule{261\,--\,0.001\,hPa} \standardrangestatement
\standardprecisionrule
\evenstatusrule
\item[Cloud consideration:] \screeningrule{Observe the higher order bits
    for the \texttt{Status} field for cloud issues, described below.}

  As an additional screen, the fifth-least-significant bit of
  \texttt{Status} (the ``low cloud'' bit) is used to flag profiles that
  may be significantly impacted by clouds.  At pressures of 147\,hPa and
  lower (higher in the atmosphere), the cloud bits may generally be
  ignored.  In the troposphere an attempt has been made to screen out
  radiances that have been influence by clouds, but some cloud-induced
  negative biases in retrieved temperature of up to 10\,K are still
  evident, particularly in the tropics.  The ``low-cloud''
  \texttt{Status} bits from the two profiles which follow a given
  profile have been found to provide significantly better screening of
  cloud-induced temperature retrieval outliers than do the profile's own
  \texttt{Status} bits.  Temperatures in the tropopause
  (261\,hPa\,--\,178\,hPa) should be rejected as possibly influenced by
  cloud if the ``low-cloud'' Status bit is set in either of the two
  profiles following the profile in question.  The screening method
  flags 16\% of tropical and 5\% of global profiles as cloudy and
  captures 86\% of the tropical 261\,hPa values for which the difference
  between MLS~T and its a~priori is more than $-$4.5\,K
  (\texttildelow2$\sigma$) below the mean of the difference.

\item[Quality field:] \qualityrule{0.65}

  The \texttt{Quality} diagnostic in \thisv is has fewer low values that
  did v2.2, reflecting better closure of the radiances used in the
  temperature retrieval.  This threshold typically excludes 1\% of
  profiles.

\item[Convergence field:] \convergencerule{1.2}

  The \texttt{Convergence} diagnostic has far fewer high values in \thisv
  than it had in v2.2, as there are far fewer poorly-converged
  ``chunks'' in the new version.  Use of this threshold typically
  discards 0.1\% of profiles, compared to 2\% or profiles flagged in
  v2.2.
\end{description}

\subsection*{Artifacts}
MLS temperature has persistent, vertically oscillating biases, in the
troposphere and stratosphere, which are believed to be due to
shortcomings in the instrument forward model and are an area of
continued research.  The impact of clouds is generally limited to
tropospheric levels in the tropics, and to a lesser extent,
mid-latitudes.  The greatest impacts of clouds are \texttildelow\,$-$10\,K, at
261\,hPa, while impacts are negligible at 100\,hPa and smaller
pressures.  Flagging of clouds is discussed above.  Biases of $>$1\,K
that were seen in v2.2, particularly in the troposphere, at the
boundaries of the nominally-10-profile ``chunks'' in which the retrieval
is processed have been greatly reduced in \thisv.  Unusually short chunks
often occur at the beginnings and ends of days and these may contain
spurious values.  Further discussion of artifacts may be found
in~\citet{SchwartzEtal08}.

\subsection*{Review of comparisons with other datasets}
\label{sec:TComparisons}

\begin{figure}
\centering\dontincludegraphics[width=0.7\linewidth]{figures/TempSummary2}
\caption{The left panel shows globally-averaged mean differences
  between MLS temperature and eight correlative data sets.  Criteria for
  coincidences are described in detain in~\citet{SchwartzEtal08}.  The
  right panel shows the global standard deviations about the means.}
\label{fig:TempSummary2}
\end{figure}

\begin{figure}
\centering\dontincludegraphics[width=0.7\linewidth]{figures/MLSv33TminusGEOS5}
\caption{Zonal mean of the difference between MLS \thisv temperature and
  GEOS-5.2 temperature (upper), and variability
  about that mean (lower), averaged for 2005--2010.}
\label{fig:MLSminusGEOS5_zonalmean}
\end{figure}

\citet{SchwartzEtal08} describes detailed comparisons of MLS v2.2
temperature with products from the Goddard Earth Observing System,
version~5~\citep{ReineckerEtal07} (GEOS-5), the European Center for
Medium-Range Weather Forecast~\citep[e.g.,][]{SimmonsEtal05} (ECMWF),
the {CHAllenging Minisatellite Payload} (CHAMP)~\citep{WickertEtal01},
the combined Atmospheric Infrared Sounder / Advanced Microwave Sounding
Unit (AIRS/AMSU), the Sounding of the Atmosphere using Broadband
Radiometry (SABER)~\citep{MlynczakAndRussell95}, the Halogen Occultation
Experiment~\citep{HervigEtal96} (HALOE) and the Atmospheric Chemistry
Experiment~\citep{BernathEtal04} (ACE), as well as to radiosondes from
the global network.  From 261\,hPa to \texttildelow10\,hPa there is generally
agreement to \texttildelow1\,K between the assimilations (ECMWF and GEOS-5) and
AIRS, radiosondes and CHAMP, with SABER and ACE having generally warm
biases of \texttildelow2\,K relative to this group.
Figure~\ref{fig:TempSummary2} shows the global mean biases in the left
panel and the 1\,$\sigma$ scatter about the mean in the right panel for
these eight comparisons.  Between 1\,hPa and 0.001\,hPa, MLS has biases
with respect to SABER of $+$1\,K~to~$-$5\,K between 1\,hPa and 0.1\,hPa,
of 0\,K~to~$-$3\,K between 0.1\,K and 0.01\,K and increasing in
magnitude to $-$10\,K at 0.001\,hPa.  Estimates of systematic error in
the MLS temperature are shown in black, with 2-$\sigma$ uncertainty
shown with gray shading.  The black line is the modeled contribution of
``gain compression,'' which was hoped would explain much of the vertical
structure of MLS biases in the upper troposphere and lower stratosphere.
As discussed above, the gain-compression model used in this study does
not adequately close the retrieval's radiance residuals, so further
study is needed to understand the forward-model inadequacies.

Figure~\ref{fig:MLSminusGEOS5_zonalmean} shows zonal mean temperature
and its variability averaged over 93 days processed with v2.2.
Persistent vertical structure in the troposphere and lower stratosphere
is evident, with the oscillations somewhat stronger at the equator and
poles than at mid-latitudes.  In the upper stratosphere, MLS has a
general warm bias relative to GEOS-5 at mid and high latitude that
increases to more than 10\,K in the poles at 1\,hPa.  The bias at 1\,hPa
is much smaller in polar summer, but persists in polar winter.

\subsection*{Desired improvements}
Improvement of the forward model, perhaps through inclusion of some
combination of amplifier non-linearity or filter shifts to
better-closed radiance residuals, would permit the concurrent use of
all of the 118-GHz and 239-GHz O$_2$ bands, and improve accuracy
throughout the profile and precision and resolution in the stratosphere. 


\begin{sidewaystable}
\caption{Summary of MLS Temperature product}\vspace{0.5ex}
\label{tab:TSummary}
\begin{minipage}{\textheight}
\centering\begin{tabular}{ccccccl}
\toprule
\textbf{Pressure} & 
\parbox[m]{4em}{\centering\textbf{Precision\footnote{Precision on individual profiles} \\/ K}} &
\parbox[m]{4em}{\centering\textbf{Observed Scatter\footnote{Precision
      inferred from differences of
      individual profiles from successive orbits (v2.2 results shown)} \\/ K}} &
\parbox[m]{4em}{\centering\textbf{Resolution\\ V $\times$ H\\ / km}} &
\parbox[m]{6em}{\centering\textbf{Modeled Bias\\uncertainty\\ / K}} &
\parbox[m]{6em}{\centering\textbf{Observed Bias\\uncertainty\\ / K}} &
\textbf{Comments} \\
\midrule
$<$0.001\,hPa      & ---       & ---      &  ---             & --- & --- & Unsuitable for scientific use \\
0.001\,hPa         & $\pm$2.5  & $\pm$3.5 & 10--13  $\times$ 220 & $+$2\,$\pm$\,3    & $-$9     & \\
0.01\,hPa          & $\pm$2.2  & $\pm$3   & 8--12  $\times$ 185 & $+$2\,$\pm$\,3    & $-$2\,to\,0 & \\
0.1\,hPa           & $\pm$2    & $\pm$2.3 & 6 $\times$ 170 & $+$2\,$\pm$\,2    & $-$8\,to\,0 & \\
0.316\,hPa         & $\pm$1    & $\pm$1.5 & 7.5 $\times$ 165 & $+$3\,$\pm$\,2    & $-$7\,to\,$-$4& \\
1\,hPa             & $\pm$1    & $\pm$1.4 & 7.0 $\times$ 165 & $+$1\,$\pm$\,2    &  0\,to\,$+$5 & \\
3.16\,hPa          & $\pm$0.8  & $\pm$1   & 5.5 $\times$ 165 & $+$2\,$\pm$\,2    & $+$1 & \\
10\,hPa            & $\pm$0.6  & $\pm$1   & 4.3 $\times$ 165 &  0\,$\pm$\,2      & $-$1\,to\,0& \\
14.7\,hPa          & $\pm$0.6  & $\pm$1   & 3.9 $\times$ 165 & $+$4\,$\pm$\,2    &  0\,to\,$+$1 & \\
31.6\,hPa          & $\pm$0.6  & $\pm$1   & 3.6 $\times$ 165 & $+$1\,$\pm$\,2    & $-$2\,to\,0 & \\
56.2\,hPa          & $\pm$0.8  & $\pm$0.8 & 3.8 $\times$ 165 & $-$1\,$\pm$\,2    & $-$2\,to\,0 & \\
100\,hPa           & $\pm$0.8  & $\pm$0.8 & 5.0 $\times$ 165 & $+$2\,$\pm$\,2    &  0\,to\,$+$1  & \\
215\,hPa           & $\pm$1    & $\pm$1   & 5.0 $\times$ 170 &  $-$1.5\,$\pm$\,4 & $-$2.5\,to\,$-$1.5   & \\
261\,hPa           & $\pm$1    & $\pm$1   & 5.0 $\times$ 170 &  0\,$\pm$\,5      & $-$2\,to\,0   & \\
1000\,--\,316\,hPa & ---       & ---      & ---              & ---           & --- & Unsuitable for scientific use\\
\bottomrule
\end{tabular}
\end{minipage}
\end{sidewaystable}

%%% Local Variables: 
%%% mode: latex
%%% TeX-master: "v4-x-report.tex"
%%% End: 
